\subsection{整数的定义}

定义1。一个基数 \(\mathfrak{a}\) 称为有限的，如果 \(\mathfrak{a} \neq \mathfrak{a} + 1\) 。有限基数也称为自然整数（如果没有混淆风险，简称为整数（*））。一个集合 \(\mathbf{E}\) 称为有限的，如果Card (E)是有限基数；此时Card (E)称为 \(\mathbf{E}\) 的元素个数.

一个族（第二章，§3，第4号）称为有限的，如果其指标集是有限的。

当我们说某种类型的对象的个数是整数m时，我们的意思是这些对象是一个有限集合的元素，该集合的元素个数为m。元素个数为m的集合也称为m个元素的集合。

命题1。基数a是有限的当且仅当a + 1是有限的。

因为关系 \(\mathfrak{a} = \mathfrak{b}\) 与 \(\mathfrak{a} + 1 = \mathfrak{b} + 1\) （基数 \(\mathfrak{a}\) 与 \(\mathfrak{b}\) 之间的关系）是等价的（§3，第4号，命题8）；因此关系 \(\mathfrak{a} \neq \mathfrak{a} + 1\) 与 \(\mathfrak{a} + 1 \neq (\mathfrak{a} + 1) + 1\) 是等价的。

显然 \(0 \neq 1\) ；因此0是整数。由此可知1和2是整数。基数 \(2 + 1\) 与 \((2 + 1) + 1\) 是整数，分别记为3和4。

